\documentclass[11pt,a4paper]{article}
\usepackage[utf8]{inputenc}
\usepackage[vietnamese]{babel}
\usepackage{amsmath,amssymb,amsfonts}
\usepackage{graphicx}
\usepackage{tikz}
\usetikzlibrary{calc,angles,quotes}
\usepackage[top=1.0cm,bottom=1.8cm,left=1.4cm,right=1.4cm,footskip=0.8cm]{geometry}
\setlength{\headheight}{16pt}
\usepackage{fancyhdr}
\usepackage{lastpage}
\usepackage{tabularx}
\usepackage{ex_test}

\pagestyle{fancy}
\fancyhf{}
\renewcommand{\headrulewidth}{0.4pt}
\renewcommand{\footrulewidth}{0.4pt}
\fancyhead[L]{\small \textbf{TOÁN 9 - KẾT NỐI TRI THỨC VỚI CUỘC SỐNG}}
\fancyhead[R]{\small \textbf{CHƯƠNG V - BÀI 17}}
\fancyfoot[L]{\small \textit{Sưu tầm \& biên soạn: Phan Đình Quân}}
\fancyfoot[R]{\small Trang \thepage/\pageref{LastPage}}

\begin{document}

% HEADER BOX
\noindent
\begin{tabularx}{\textwidth}{|>{\centering\arraybackslash}p{6.8cm}|>{\centering\arraybackslash}X|}
\hline
\textbf{CHƯƠNG V: ĐƯỜNG TRÒN} & \textbf{PHIẾU BÀI TẬP TRẮC NGHIỆM} \\
\textbf{BÀI 17: VỊ TRÍ TƯƠNG ĐỐI CỦA HAI ĐƯỜNG TRÒN} & \textbf{Môn: TOÁN 9 (Bộ sách Kết nối tri thức)} \\
\textit{Thời gian làm bài: 45 phút} & \textit{Họ và tên: ...........................................................} \\
\hline
\end{tabularx}

\vspace{0.3cm}

\begin{flushleft}
	{\color{blue!50!green}\fbox{\fontfamily{qag}\bfseries\selectfont PHẦN 1. Câu trắc nghiệm 4 phương án}}
\end{flushleft}
\setcounter{ex}{0}

\begin{ex}
Nếu hai đường tròn không cắt nhau thì số điểm chung của hai đường tròn là
\choice
{$1$}
{$2$}
{$3$}
{\True $0$}
\loigiai{}
\end{ex}

\begin{ex}
Cho hai đường tròn $(O_1)$ và $(O_2)$ tiếp xúc ngoài tại $A$ và một đường thẳng $\mathrm{d}$ tiếp xúc với $(O_1)$ và $(O_2)$ lần lượt tại $B, C$. Tam giác $ABC$ là
\choice
{Tam giác cân}
{Tam giác đều}
{\True Tam giác vuông}
{Tam giác vuông cân}
\loigiai{}
\end{ex}

\begin{ex}
Cho đoạn thẳng $OO'$ và điểm $A$ nằm trên đoạn $OO'$ sao cho $OA = 2O'A$. Vị trí tương đối của đường tròn tâm $O$ bán kính $OA$ và đường tròn tâm $O'$ bán kính $O'A$ là
\choice
{Nằm ngoài nhau}
{Cắt nhau}
{\True Tiếp xúc ngoài}
{Tiếp xúc trong}
\loigiai{}
\end{ex}

\begin{ex}
Cho hai đường tròn $(O), (O')$ cắt nhau tại $A, B$ trong đó $O' \in (O)$. Kẻ đường kính $O'C$ của đường tròn $(O)$. Chọn khẳng định \textbf{sai}.
\choice
{$AC = CB$}
{$\widehat{CBO'} = 90^\circ$}
{$CA, CB$ là hai tiếp tuyến của $(O')$}
{\True $CA, CB$ là hai cát tuyến của $(O')$}
\loigiai{}
\end{ex}

\begin{ex}
Cho hai đường tròn tiếp xúc ngoài $(O; R)$ và $(O'; r)$ với $R > r$ và khoảng cách giữa hai tâm $OO' = \mathrm{d}$. Khi đó
\choice
{$\mathrm{d} = R - r$}
{$\mathrm{d} > R + r$}
{$R - r < \mathrm{d} < R + r$}
{\True $\mathrm{d} = R + r$}
\loigiai{}
\end{ex}

\begin{ex}
Cho đường tròn $(O_1; 3\text{ cm})$ tiếp xúc ngoài với đường tròn $(O_2; 1\text{ cm})$. Vẽ bán kính $O_1B$ và $O_2C$ song song với nhau cùng thuộc nửa mặt phẳng bờ $O_1O_2$. Gọi $D$ là giao điểm của $BC$ và $O_1O_2$. Tính số đo góc $\widehat{BAC}$ (với $A$ là tiếp điểm của $(O_1)$ và $(O_2)$).
\choice
{\True $90^\circ$}
{$60^\circ$}
{$100^\circ$}
{$80^\circ$}
\loigiai{}
\end{ex}

\begin{ex}
Cho hai đường tròn $(O; 10\text{ cm})$ và $(O'; 5\text{ cm})$ cắt nhau tại $A, B$. Tính độ dài đoạn thẳng $OO'$ biết $AB = 8\text{ cm}$ và $O, O'$ nằm cùng phía đối với $AB$ (làm tròn kết quả đến chữ số thập phân thứ nhất).
\choice
{$OO' \approx 6{,}5\text{ cm}$}
{$OO' \approx 6{,}1\text{ cm}$}
{$OO' \approx 6{,}3\text{ cm}$}
{\True $OO' \approx 6{,}2\text{ cm}$}
\loigiai{}
\end{ex}

\begin{ex}
Cho hai đường tròn $(O; 20\text{ cm})$ và $(O'; 15\text{ cm})$ cắt nhau tại $A, B$. Tính độ dài đoạn thẳng $OO'$ biết $AB = 24\text{ cm}$ và $O, O'$ nằm cùng phía đối với $AB$.
\choice
{\True $OO' = 7\text{ cm}$}
{$OO' = 8\text{ cm}$}
{$OO' = 9\text{ cm}$}
{$OO' = 25\text{ cm}$}
\loigiai{}
\end{ex}

\newpage

\begin{ex}
Cho hai đường tròn $(O), (O')$ tiếp xúc ngoài tại $A$. Kẻ tiếp tuyến chung ngoài $MN$ với $M \in (O), N \in (O')$. Gọi $P, Q$ lần lượt là điểm đối xứng với $M, N$ qua đường nối tâm $OO'$. Khi đó tứ giác $MNPQ$ là hình gì?
\choice
{\True Hình thang cân}
{Hình thang}
{Hình thang vuông}
{Hình bình hành}
\loigiai{}
\end{ex}

\begin{ex}
Cho đường tròn $(O; 6\text{ cm})$ và $(O'; 2\text{ cm})$ cắt nhau tại $A, B$ sao cho $OA$ là tiếp tuyến của $(O')$. Độ dài dây chung $AB$ bằng
\choice
{$AB = 3\sqrt{10}\text{ cm}$}
{\True $AB = \dfrac{6\sqrt{10}}{5}\text{ cm}$}
{$AB = \dfrac{3\sqrt{10}}{5}\text{ cm}$}
{$AB = \dfrac{\sqrt{10}}{5}\text{ cm}$}
\loigiai{}
\end{ex}

\begin{ex}
Cho các đường tròn $(O_1; 10\text{ cm})$, $(O_2; 15\text{ cm})$ và $(O_3; 15\text{ cm})$ tiếp xúc ngoài với nhau đôi một. Hai đường tròn $(O_2)$ và $(O_3)$ tiếp xúc nhau tại điểm $A$. Đường tròn $(O_1)$ tiếp xúc với $(O_2)$ và $(O_3)$ lần lượt tại $C$ và $B$. Khẳng định nào sau đây là đúng?
\choice
{\True $O_1A$ là tiếp tuyến chung của đường tròn $(O_2)$ và $(O_3)$}
{$O_1A = 25\text{ cm}$}
{$O_1A = 15\text{ cm}$}
{Cả A và B đều đúng}
\loigiai{}
\end{ex}

\begin{ex}
Cho hai đường tròn $(O), (O')$ tiếp xúc ngoài tại $A$. Kẻ các đường kính $AB$ của $(O)$ và $AC$ của $(O')$, gọi $DE$ là tiếp tuyến chung ngoài của hai đường tròn $(D \in (O), E \in (O'))$. Gọi $M$ là giao điểm của $BD$ và $CE$. Tính diện tích tứ giác $ADME$ biết $\widehat{DOA} = 60^\circ$ và $OA = 6\text{ cm}$.
\choice
{\True $12\sqrt{3}\text{ cm}^2$}
{$12\text{ cm}^2$}
{$16\text{ cm}^2$}
{$24\text{ cm}^2$}
\loigiai{}
\end{ex}

\begin{ex}
Cho hai đường tròn $(O), (O')$ tiếp xúc ngoài tại $A$. Kẻ các đường kính $AB$ của $(O)$ và $AC$ của $(O')$, gọi $DE$ là tiếp tuyến chung ngoài của hai đường tròn $(D \in (O), E \in (O'))$. Gọi $M$ là giao điểm của $BD$ và $CE$. Tính diện tích tứ giác $ADME$ biết $\widehat{DOA} = 60^\circ$ và $OA = 8\text{ cm}$.
\choice
{$12\sqrt{3}\text{ cm}^2$}
{\True $\dfrac{64\sqrt{3}}{3}\text{ cm}^2$}
{$\dfrac{32\sqrt{3}}{3}\text{ cm}^2$}
{$36\text{ cm}^2$}
\loigiai{}
\end{ex}

\begin{ex}
Cho các đường tròn $(O_1; 10\text{ cm})$, $(O_2; 15\text{ cm})$ và $(O_3; 15\text{ cm})$ tiếp xúc ngoài với nhau đôi một. Hai đường tròn $(O_2)$ và $(O_3)$ tiếp xúc nhau tại điểm $A$. Đường tròn $(O_1)$ tiếp xúc với $(O_2)$ và $(O_3)$ lần lượt tại $C$ và $B$. Tính diện tích tam giác $ABC$.
\choice
{$36\text{ cm}^2$}
{\True $72\text{ cm}^2$}
{$144\text{ cm}^2$}
{$96\text{ cm}^2$}
\loigiai{}
\end{ex}

\begin{ex}
Cho hai đường tròn $(O; R)$ và $(O'; r)$ với $R > r$ tiếp xúc ngoài tại $A$. Vẽ các bán kính $OB \parallel O'D$ với $B, D$ cùng thuộc nửa mặt phẳng bờ $OO'$. Đường thẳng $DB$ và $OO'$ cắt nhau tại $I$. Tiếp tuyến chung ngoài $GH$ của $(O)$ và $(O')$ với $G, H$ nằm ở nửa mặt phẳng bờ $OO'$ không chứa $B, D$. Tính $OI$ theo $R$ và $r$.
\choice
{$OI = \dfrac{R+r}{R-r}$}
{$OI = \dfrac{R-r}{R+r}$}
{$OI = \dfrac{R(R-r)}{R+r}$}
{\True $OI = \dfrac{R(R+r)}{R-r}$}
\loigiai{}
\end{ex}

\vspace{0.4cm}
\begin{center}
\textbf{BẢNG ĐÁP ÁN TRẮC NGHIỆM} \\[0.2cm]
\begin{tabular}{|c|c|c|c|c|c|c|c|c|c|c|c|c|c|c|}
\hline
\textbf{1} & \textbf{2} & \textbf{3} & \textbf{4} & \textbf{5} & \textbf{6} & \textbf{7} & \textbf{8} & \textbf{9} & \textbf{10} & \textbf{11} & \textbf{12} & \textbf{13} & \textbf{14} & \textbf{15} \\
\hline
\textbf{D} & \textbf{C} & \textbf{C} & \textbf{D} & \textbf{D} & \textbf{A} & \textbf{D} & \textbf{A} & \textbf{A} & \textbf{B} & \textbf{A} & \textbf{A} & \textbf{B} & \textbf{B} & \textbf{D} \\
\hline
\end{tabular}
\end{center}

\vspace{0.3cm}
\begin{center}
\textbf{--- HẾT ---}
\end{center}

\end{document}
